\documentclass{ltjsarticle}

\usepackage{graphicx}
\usepackage{amsmath}
\usepackage{comment}
\usepackage{color}
\usepackage{url}
\usepackage{siunitx}
\usepackage[version=4]{mhchem}
\usepackage{paralist}
\usepackage{longtable}
\usepackage{multirow}

\usepackage{luatexja}

\usepackage{float}
\usepackage{minted}
\usepackage{tikz}
\usepackage{circuitikz}
\usepackage{bm}
\usepackage[superscript]{cite}
\usepackage{physics}
\usepackage{subcaption}
\usepackage{placeins}
\usepackage{wrapfig}

\usepackage{newfloat}

\DeclareFloatingEnvironment[
  fileext=lol,
  listname=List of Cord Listings,
  name=Listing
]{cord_listing}

\usepackage{enumitem}
\setlist[description]{parsep=5pt}
\setlist[enumerate]{parsep=5pt}

\usepackage{hyperref}
\usepackage[nameinlink]{cleveref}



\begin{document}

\section{6/4: Gradle+Kotlin のセットアップ}
\begin{minted}[frame=single]{bash}
    gradle --version
\end{minted}
でversion の確認を行った，8.x.x だったため，
\begin{minted}[frame=single]{bash}
    gradle init
    ./gradlew wrapper --gradle-version 9.5.1   
\end{minted}
を行った．
\begin{minted}[frame=single]{bash}
    ./gradlew run    
\end{minted}
で\texttt{Hello World!}の出力が得られた．これは，app/src/main/kotlin/org/example/App.kt 内の\texttt{main()}関数が実行されたものである．
実行される関数は，app/src/build.gradle.kts 内の
\begin{minted}[frame=single]{kotlin}
    application {
        mainClass = "org.example.AppKt"
    }    
\end{minted}
の設定によって決まった．

.../org/example/SerializationDemo.kt を作成して，ライブラリの有効化を試した．
gradle/libs.versions.toml の\texttt{[plugins]}に追加して，build.gradle.kts の\texttt{plugins}に\texttt{alias()}で追加する．
同じく\texttt{[libraries]}に追加して，build.gradle.kts の\texttt{dependencies}に\texttt{implementation()}で追加した．
\end{document}
